\section{2024 Geometry \& Topology}

\subsection{Exam}

\begin{exercise}
\label{geometry:2024:1}
Let $n > 1$ be a positive integer.
(i) Does there exist a map $f : S^{2n} \to \mathbb{CP}^{n}$ with $\deg(f) \neq 0$? Construct an example or disprove it.
(ii) Does there exist a map $f : \mathbb{CP}^{n} \to S^{2n}$ with $\deg(f) \neq 0$? Construct an example or disprove it.
\end{exercise}

\begin{solution}
\textbf{Prerequisites \& Definitions:}
\begin{itemize}
    \item \textbf{Complex Projective Space $\mathbb{CP}^n$}: This is the space of all 1-dimensional complex subspaces (lines through the origin) in $\mathbb{C}^{n+1}$. As a manifold, its real dimension is $2n$\footnote{To see the dimension, note that a point in $\mathbb{CP}^n$ is represented by a non-zero vector $(z_0, \dots, z_n) \in \mathbb{C}^{n+1}$, with the identification $(z_0, \dots, z_n) \sim \lambda(z_0, \dots, z_n)$ for $\lambda \in \mathbb{C} \setminus \{0\}$. We can normalize so that $\sum |z_i|^2 = 1$, which gives $S^{2n+1}$, and then quotient by the action of $S^1$ (the unit complex numbers), so $\dim_{\mathbb{R}}(\mathbb{CP}^n) = (2n+1) - 1 = 2n$.}.
    \item \textbf{Cohomology Ring}: For $\mathbb{CP}^n$, the integral cohomology is $H^*(\mathbb{CP}^n; \mathbb{Z}) \cong \mathbb{Z}[x]/(x^{n+1})$, where $x \in H^2(\mathbb{CP}^n; \mathbb{Z})$ is the generator\footnote{The generator $x$ is the first Chern class $c_1(\mathcal{L})$ of the tautological line bundle $\mathcal{L} \to \mathbb{CP}^n$. Geometrically, it represents the dual of the hyperplane $\mathbb{CP}^{n-1} \subset \mathbb{CP}^n$.}.
    \item \textbf{Degree of a Map}: For a map $f: M \to N$ between oriented closed manifolds of the same dimension $d$, the degree $\deg(f)$ is the integer such that $f^*([N]) = \deg(f)[M]$, where $[M] \in H^d(M)$ is the fundamental class\footnote{Intuitively, the degree counts how many times the domain "wraps around" the target, accounting for orientation.}.
\end{itemize}

\textbf{Geometric Intuition:}
The "richness" of the cohomology ring tells us about the topological structure. $\mathbb{CP}^n$ has non-trivial topology in every even dimension $2, 4, \dots, 2n$, tied together by the relation $x^k = x \cup \dots \cup x$. In contrast, the sphere $S^{2n}$ is topologically "hollow" — it has no non-trivial cohomology between dimensions $0$ and $2n$. Part (i) asks if a "hollow" sphere can map onto the "structured" $\mathbb{CP}^n$ while preserving the top-dimensional volume (degree). The ring structure says "no": the top class of $\mathbb{CP}^n$ is built from the second-dimensional class, which doesn't exist on the sphere. Part (ii) asks the reverse: can we "collapse" the internal structure of $\mathbb{CP}^n$ to just look like a sphere at the top? The answer is "yes", because we can just pinch everything below the top dimension to a point.

\textbf{Formal Proof:}

\textbf{(i) Non-existence of non-zero degree map $f: S^{2n} \to \mathbb{CP}^n$ for $n > 1$.}
Suppose such a map $f$ exists. Consider the pullback homomorphism on cohomology $f^*: H^*(\mathbb{CP}^n; \mathbb{Z}) \to H^*(S^{2n}; \mathbb{Z})$.
For $n > 1$, the dimension of the sphere is $2n \ge 4$. The second cohomology group of the sphere is:
\sidenote{$H^2(S^{2n}; \mathbb{Z}) = 0$ for $2n > 2$.}
Let $x \in H^2(\mathbb{CP}^n; \mathbb{Z})$ be the generator. Since $f^*$ is a ring homomorphism, we have:
\sidenote{$f^*(x^n) = (f^*x)^n$.}
However, $f^*x$ must lie in $H^2(S^{2n}; \mathbb{Z})$, which is $0$. Thus $f^*x = 0$.
It follows that $f^*(x^n) = 0^n = 0$.
By the definition of degree:
\sidenote{$f^*(x^n) = \deg(f) \cdot [S^{2n}]$.}
Since $x^n$ is the generator of $H^{2n}(\mathbb{CP}^n; \mathbb{Z})$ and $f^*(x^n) = 0$, we must have $\deg(f) = 0$.
Thus, no map of non-zero degree exists.

\textbf{(ii) Existence of a map $f: \mathbb{CP}^n \to S^{2n}$ with $\deg(f) = 1$.}
We can construct such a map using the CW complex structure of $\mathbb{CP}^n$.
Recall that $\mathbb{CP}^n$ can be constructed by attaching cells of dimensions $0, 2, 4, \dots, 2n$\footnote{Specifically, $\mathbb{CP}^n$ is obtained from $\mathbb{CP}^{n-1}$ by attaching a $2n$-cell $e^{2n}$ via the Hopf map $S^{2n-1} \to \mathbb{CP}^{n-1}$.}.
The $(2n-2)$-skeleton of $\mathbb{CP}^n$ is exactly $\mathbb{CP}^{n-1}$.
Consider the quotient space $\mathbb{CP}^n / \mathbb{CP}^{n-1}$.
\sidenote{$\mathbb{CP}^n / \mathbb{CP}^{n-1} \cong S^{2n}$.}
Let $q: \mathbb{CP}^n \to \mathbb{CP}^n / \mathbb{CP}^{n-1} \cong S^{2n}$ be the quotient map that collapses the entire $(2n-2)$-skeleton to a single point.
This map $q$ is a homeomorphism on the interior of the top $2n$-cell and maps its boundary to the collapsed point.
In cohomology, the induced map $q^*: H^{2n}(S^{2n}; \mathbb{Z}) \to H^{2n}(\mathbb{CP}^n; \mathbb{Z})$ is an isomorphism\footnote{This is because $H^{2n}(\mathbb{CP}^n, \mathbb{CP}^{n-1}; \mathbb{Z}) \cong \tilde{H}^{2n}(\mathbb{CP}^n / \mathbb{CP}^{n-1}; \mathbb{Z}) \cong H^{2n}(S^{2n}; \mathbb{Z})$, and the long exact sequence of the pair $(\mathbb{CP}^n, \mathbb{CP}^{n-1})$ shows this isomorphism maps to the generator of $H^{2n}(\mathbb{CP}^n)$.}.
Thus, $\deg(q) = 1$.

\begin{remark}
For $n \ge 1$, the integral homology and cohomology groups of the $n$-dimensional sphere $S^n$ and the complex projective space $\mathbb{CP}^n$ are:
\begin{itemize}
    \item \textbf{Sphere $S^n$}:
    \begin{equation*}
        H_k(S^n; \mathbb{Z}) \cong \begin{cases} \mathbb{Z} & k=0, n \\ 0 & \text{otherwise} \end{cases}, \quad H^k(S^n; \mathbb{Z}) \cong \begin{cases} \mathbb{Z} & k=0, n \\ 0 & \text{otherwise} \end{cases}
    \end{equation*}
    \item \textbf{Complex Projective Space $\mathbb{CP}^n$}:
    \begin{equation*}
        H_k(\mathbb{CP}^n; \mathbb{Z}) \cong \begin{cases} \mathbb{Z} & k \in \{0, 2, \dots, 2n\} \\ 0 & \text{otherwise} \end{cases}, \quad H^k(\mathbb{CP}^n; \mathbb{Z}) \cong \begin{cases} \mathbb{Z} & k \in \{0, 2, \dots, 2n\} \\ 0 & \text{otherwise} \end{cases}
    \end{equation*}
\end{itemize}
\footnote{The cohomology ring for $S^n$ is $H^*(S^n; \mathbb{Z}) \cong \mathbb{Z}[\alpha]/(\alpha^2)$ where $|\alpha| = n$. For $\mathbb{CP}^n$, it is $H^*(\mathbb{CP}^n; \mathbb{Z}) \cong \mathbb{Z}[x]/(x^{n+1})$ where $|x| = 2$. For $n > 1$, this implies $H^2(S^{2n}; \mathbb{Z}) = 0$ while $H^2(\mathbb{CP}^n; \mathbb{Z}) \cong \mathbb{Z}$. This difference is the critical fact used to show $\deg(f)=0$ in part (i), as any map $f: S^{2n} \to \mathbb{CP}^n$ must pull back the generator $x \in H^2(\mathbb{CP}^n)$ to zero, forcing $f^*(x^n) = (f^*x)^n = 0$.}
\end{remark}
\end{solution}

\begin{exercise}
\label{geometry:2024:2}
Let $\Sigma \subset \mathbb{R}^3$ be an embedded surface in $\mathbb{R}^3$. A surface is called minimal if, for any $p \in \Sigma$, we have $\kappa_1(p) + \kappa_2(p) = 0$, where $\kappa_1(p)$ and $\kappa_2(p)$ are the two principal curvatures at $p$. Prove that if $\Sigma$ is closed, then $\Sigma$ cannot be minimal.
\end{exercise}

\begin{solution}
\textbf{Prerequisites \& Definitions:}
\begin{itemize}
    \item \textbf{Principal Curvatures $\kappa_1, \kappa_2$}: At any point $p$ on a surface $\Sigma$, these are the eigenvalues of the shape operator $S_p = -d\mathbf{N}_p$\footnote{Geometrically, they represent the maximum and minimum normal curvatures of all curves on the surface passing through $p$. If we rotate a normal plane around the normal vector, the curvature of the intersection curve varies between these two values.}.
    \item \textbf{Mean Curvature $H$}: Defined as $H = \frac{\kappa_1 + \kappa_2}{2}$. A surface is \textbf{minimal} if $H = 0$ everywhere\footnote{The term "minimal" comes from the fact that these surfaces locally minimize area. Specifically, they are critical points for the area functional under compactly supported variations.}.
    \item \textbf{Closed Surface}: A surface $\Sigma$ is called \textbf{closed} if it is a compact manifold without boundary\footnote{In the context of $\mathbb{R}^3$, being closed means the surface is bounded and "loops back" on itself, having no edges or boundaries (like a sphere or a torus). This topological constraint is crucial because it implies the existence of extreme points (maxima/minima of coordinate functions) and allows for the application of the Maximum Principle.}.
    \item \textbf{Properties of a Closed Surface in $\mathbb{R}^3$}:
    \begin{itemize}
        \item \textbf{Compactness}: Since $\Sigma$ is closed and bounded in $\mathbb{R}^3$, it is compact by the Heine-Borel theorem.
        \item \textbf{Orientability}: A fundamental result in topology states that any closed, embedded surface in $\mathbb{R}^3$ is orientable.
        \item \textbf{Jordan-Brouwer Separation Theorem}: A closed surface in $\mathbb{R}^3$ divides the space into exactly two regions: a bounded "inside" and an unbounded "outside."
    \end{itemize}
    \item \textbf{Laplace-Beltrami Operator $\Delta_{\Sigma}$}: For a surface $\Sigma$, the Laplacian of the position vector $\vec{x}$ is given by $\Delta_{\Sigma} \vec{x} = 2H\mathbf{N}$. For a minimal surface ($H=0$), the coordinate functions are harmonic.
\end{itemize}

\textbf{Geometric Intuition:}
The condition $\kappa_1 + \kappa_2 = 0$ implies that at every point, the surface must look like a "saddle" (unless it is flat, $\kappa_1 = \kappa_2 = 0$). If it curves "up" in one direction, it must curve "down" in the perpendicular direction.
However, a closed surface in $\mathbb{R}^3$ is "trapped" within a finite region. If you look at the "highest" point (the one with the largest $z$-coordinate), the surface must curve "down" in \textit{all} directions to stay below that maximum level. This requires both $\kappa_1$ and $\kappa_2$ to be non-positive ($\kappa_1, \kappa_2 \le 0$). But if their sum is zero, then they both must be zero. This suggests the surface must be flat at its extreme points, which is a very strong constraint that ultimately leads to a contradiction when combined with the global structure.

\textbf{Formal Proof:}

\begin{lemma}[Maximum Principle for Harmonic Functions]
Let $M$ be a connected, compact Riemannian manifold without boundary. If $u \in C^2(M)$ is a harmonic function (i.e., $\Delta_M u = 0$), then $u$ must be constant.
\end{lemma}

\begin{proof}
Consider the integral of the Laplacian of $u^2$ over the manifold $M$. By the identity $\Delta(u^2) = 2u\Delta u + 2|\nabla u|^2$, and the fact that $\Delta u = 0$, we have:
\[
\Delta(u^2) = 2|\nabla u|^2.
\]
Integrating this over the entire manifold $M$:
\[
\int_M \Delta(u^2) \, dV_g = 2 \int_M |\nabla u|^2 \, dV_g.
\]
By the Divergence Theorem, since $M$ is a compact manifold without boundary ($\partial M = \emptyset$):
\sidenote{$\int_M \Delta f \, dV_g = \int_{\partial M} \frac{\partial f}{\partial \nu} \, dA = 0$.}
Thus, the integral of the Laplacian must be zero:
\[
0 = 2 \int_M |\nabla u|^2 \, dV_g.
\]
Since $|\nabla u|^2 \ge 0$, the integrand must vanish identically on $M$. This implies $\nabla u = 0$ everywhere on $M$. Because $M$ is connected, a function with a zero gradient must be constant.
\end{proof}

\textbf{Method 1: The Harmonic Function Argument (Analytical)}
Assume $\Sigma \subset \mathbb{R}^3$ is a minimal, closed, embedded surface.
By definition, the mean curvature $H$ is identically zero on $\Sigma$.
Recall the fundamental formula for the Laplace-Beltrami operator acting on the position vector $\vec{x} = (x_1, x_2, x_3)$ of the surface:
\sidenote{Derivation of $\Delta_{\Sigma} \vec{x} = 2H\mathbf{N}$:
In local coordinates $\{u^i\}$, the Laplacian is $\Delta f = g^{ij} (\partial_i \partial_j f - \Gamma_{ij}^k \partial_k f)$.
Applying this to $\vec{x}$, we have $\Delta_{\Sigma} \vec{x} = g^{ij} (\vec{x}_{ij} - \Gamma_{ij}^k \vec{x}_k)$.
From the Gauss formula, $\vec{x}_{ij} = \Gamma_{ij}^k \vec{x}_k + h_{ij} \mathbf{N}$, where $h_{ij}$ is the second fundamental form.
The tangential terms cancel: $\Delta_{\Sigma} \vec{x} = g^{ij} h_{ij} \mathbf{N}$.
By definition, $H = \frac{1}{2} \text{tr}(g^{-1} h) = \frac{1}{2} g^{ij} h_{ij}$, thus $\Delta_{\Sigma} \vec{x} = 2H\mathbf{N}$.}
Since $H = 0$, we have $\Delta_{\Sigma} \vec{x} = 0$. This implies that each coordinate function $x_i : \Sigma \to \mathbb{R}$ is a \textbf{harmonic function} on the manifold $\Sigma$:
\sidenote{$\Delta_{\Sigma} x_i = 0$ for $i=1, 2, 3$.}
Now, we utilize the \textbf{Maximum Principle} for harmonic functions on a compact manifold without boundary.
\footnote{The Maximum Principle states that a harmonic function $u$ on a connected compact Riemannian manifold without boundary must be constant. If $u$ were not constant, it would have a maximum at some point $p$. At a maximum, the Hessian $\nabla^2 u$ must be negative semi-definite, which implies $\Delta u = \text{tr}(\nabla^2 u) \le 0$. If $\Delta u = 0$, the function must be "flat" in a way that forces it to be constant by the strong maximum principle.}
Since $\Sigma$ is closed (compact and without boundary), and each $x_i$ is harmonic on $\Sigma$, each $x_i$ must be constant on each connected component of $\Sigma$.
\sidenote{$x_1 \equiv c_1, \ x_2 \equiv c_2, \ x_3 \equiv c_3$ for some constants $c_i$.}
This implies that the position vector $\vec{x}$ is constant across each connected component of $\Sigma$. Consequently, the entire component is mapped to a single point $(c_1, c_2, c_3)$ in $\mathbb{R}^3$.
\sidenote{Dimension Contradiction:
An embedded surface $\Sigma \subset \mathbb{R}^3$ must locally be homeomorphic to $\mathbb{R}^2$ and thus have dimension 2. A single point has dimension 0. This topological collapse contradicts the assumption that $\Sigma$ is a manifold of dimension 2.}
Therefore, a closed minimal surface cannot be embedded in $\mathbb{R}^3$.

\textbf{Method 2: The Extreme Point Argument (Geometric)}
Consider the function $f(\vec{x}) = |\vec{x}|^2 = x_1^2 + x_2^2 + x_3^2$ defined on $\Sigma$.
Since $\Sigma$ is compact, $f$ must attain a maximum at some point $p_0 \in \Sigma$.
Let $R^2 = f(p_0)$ be the maximum value. Then the entire surface $\Sigma$ is contained within the closed ball $B_R(0)$.
At the point $p_0$, the surface $\Sigma$ is tangent to the sphere $S_R(0)$.
\sidenote{Curvature Comparison:
At the point $p_0$, the principal curvatures of the surface $\Sigma$ (relative to the inward normal) must be at least as large as the principal curvatures of the enclosing sphere $S_R(0)$.}
The principal curvatures of the sphere $S_R(0)$ are both $1/R$.
Thus, for the surface $\Sigma$ at $p_0$, we must have $\kappa_1 \ge 1/R$ and $\kappa_2 \ge 1/R$.
This implies that the mean curvature at $p_0$ satisfies:
\sidenote{$H(p_0) = \frac{\kappa_1 + \kappa_2}{2} \ge \frac{1}{R} + \frac{1}{R} \cdot \frac{1}{2} = \frac{1}{R} > 0$.}
This contradicts the assumption that $\Sigma$ is a minimal surface ($H \equiv 0$).
Thus, $\Sigma$ cannot be minimal.
\end{solution}

\begin{remark}[The Landscape of Surface Curvatures]
The local and global geometry of a surface in $\mathbb{R}^3$ is characterized by a hierarchy of curvature invariants, from local principal curvatures to high-dimensional tensors:
\begin{enumerate}
    \item \textbf{Principal Curvatures ($\kappa_1, \kappa_2$)}: These are the eigenvalues of the shape operator $S = -d\mathbf{N}$. 
    \sidenote{The Weingarten Equation:
    The shape operator (or Weingarten map) $S_p: T_p \Sigma \to T_p \Sigma$ is defined by $S_p(v) = -\nabla_v \mathbf{N}$. In local coordinates, this gives the Weingarten equations: $\partial_j \mathbf{N} = -S^i_j \partial_i$, where $S^i_j = g^{ik} h_{kj}$.}
    \item \textbf{Gauss Curvature ($K$)}: Defined extrinsically as the product of principal curvatures ($K = \kappa_1 \kappa_2 = \det S$). 
    \sidenote{Theorema Egregium:
    Gauss's "Remarkable Theorem" states that $K$ is actually an \textbf{intrinsic} invariant. It depends only on the metric $g_{ij}$ and its derivatives, and is related to the Riemann tensor by the Gauss equation: $R_{1212} = K \det(g)$.}
    \item \textbf{Mean Curvature ($H$)}: Defined as the average curvature $H = \frac{1}{2}(\kappa_1 + \kappa_2) = \frac{1}{2} \text{tr}(S)$. It is an \textbf{extrinsic} measure of how the surface is "bent" in $\mathbb{R}^3$.
    \item \textbf{Riemann Curvature Tensor ($R$)}: A $(1,3)$-tensor $R(X, Y)Z = \nabla_X \nabla_Y Z - \nabla_Y \nabla_X Z - \nabla_{[X, Y]} Z$ that measures the failure of covariant derivatives to commute. 
    \sidenote{Riemann Tensor in 2D:
    For a 2-manifold, the tensor is entirely determined by $K$. In its $(0,4)$-form, $R_{ijkl} = K(g_{ik}g_{jl} - g_{il}g_{jk})$. This form satisfies all symmetry properties of the Riemann tensor (skew-symmetry in $ij$ and $kl$, and symmetry under $ij \leftrightarrow kl$).}
    \item \textbf{Sectional Curvature ($Sec$)}: A generalization of Gauss curvature to higher dimensions. For a plane $P \subset T_p M$ spanned by orthonormal vectors $\{u, v\}$, $Sec(P) = \langle R(u, v)v, u \rangle$. For surfaces, $Sec(T_p \Sigma) = K$.
    \item \textbf{Ricci Curvature ($\mathrm{Ric}$)}: A $(0,2)$-tensor defined by $\mathrm{Ric}(V, W) = \sum_{i=1}^n \langle R(e_i, V)W, e_i \rangle$. 
    \sidenote{Ricci Tensor in 2D:
    For a surface ($n=2$), the Ricci tensor is isotropic: $\mathrm{Ric} = K g$. This reflects the fact that there is only one independent sectional curvature at each point.}
    \item \textbf{Scalar Curvature ($R_{scal}$)}: The trace of the Ricci tensor $R_{scal} = g^{ij} R_{ij}$. In 2D, $R_{scal} = 2K$. 
    \sidenote{Convention Note:
    In Exercise \ref{geometry:2024:4}, the symbol $S_p$ is defined as $\frac{1}{n} R_{scal}$. For a surface, this convention yields $S_p = \frac{1}{2}(2K) = K$.}
\end{enumerate}

\textbf{Classification of Points on a Surface}:
The signs of $K$ and $H$ classify the local shape at point $p$:
\begin{itemize}
    \item \textbf{Elliptic} ($K > 0$): Same signs ($\kappa_1 \kappa_2 > 0$). Locally convex (e.g., sphere).
    \item \textbf{Hyperbolic} ($K < 0$): Opposite signs ($\kappa_1 \kappa_2 < 0$). Locally saddle-shaped (characteristic of minimal surfaces).
    \item \textbf{Parabolic} ($K = 0, H \neq 0$): One zero curvature. Locally cylindrical.
    \item \textbf{Planar} ($K = 0, H = 0$): Both zero. Locally flat.
    \item \textbf{Umbilic} ($\kappa_1 = \kappa_2$): Curvature isotropic in all directions ($H^2 = K$).
\end{itemize}

\textbf{Geometric Identities}:
The principal curvatures are the roots of $\lambda^2 - 2H\lambda + K = 0$:
\[ \kappa_{1,2} = H \pm \sqrt{H^2 - K}. \]
For a minimal surface ($H=0$), we have $\kappa_{1,2} = \pm \sqrt{-K}$, implying $K \le 0$ everywhere, with $K < 0$ at non-planar points.
\end{remark}

\begin{exercise}
\label{geometry:2024:3}
Let $M$ be a closed, simply connected 6-dimensional manifold. Suppose $H_2(M) = \mathbb{Z}_2$. Prove that the Euler characteristic $\chi(M) \neq -1$.
\end{exercise}

\begin{solution}
\textbf{Prerequisites \& Definitions:}
\begin{itemize}
    \item \textbf{Simply Connected}: A manifold $M$ is simply connected if its fundamental group $\pi_1(M)$ is trivial. This implies that the first Betti number $b_1 = 0$ and $H_1(M; \mathbb{Z}) = 0$\footnote{The Hurewicz Theorem states that for a path-connected space, the first non-zero homotopy group and homology group occur at the same dimension and are isomorphic. Thus, $\pi_1(M) = 0 \implies H_1(M; \mathbb{Z}) = 0$.}.
    \item \textbf{Poincaré Duality}: For a closed oriented $n$-manifold, there is an isomorphism $H^k(M; \mathbb{Z}) \cong H_{n-k}(M; \mathbb{Z})$. This implies that the Betti numbers satisfy the symmetry $b_k = b_{n-k}$\footnote{Poincaré Duality is a deep result reflecting the self-duality of the manifold's structure. It ensures that the "holes" in dimension $k$ are perfectly mirrored by "holes" in dimension $n-k$.}.
    \item \textbf{Universal Coefficient Theorem (UCT)}: Relates cohomology to homology. Specifically, $H^k(M; \mathbb{Z}) \cong \text{Hom}(H_k(M), \mathbb{Z}) \oplus \text{Ext}(H_{k-1}(M), \mathbb{Z})$. This theorem allows us to track how torsion in homology affects cohomology.
    \item \textbf{Intersection Form}: For a $2n$-dimensional manifold, the cup product followed by evaluation on the fundamental class $[M]$ defines a bilinear form on $H^n(M)$.
    \item \textbf{Symplectic Vector Space}: A finite-dimensional vector space $V$ over a field (typically $\mathbb{R}$) equipped with a \textbf{symplectic form} $\omega$, which is a non-degenerate, skew-symmetric bilinear form $\omega: V \times V \to \mathbb{R}$\footnote{A fundamental property of symplectic vector spaces is that they must be even-dimensional. This is because a non-degenerate skew-symmetric matrix $J$ must have $\det(J) \neq 0$, and since $J^T = -J$, we have $\det(J) = \det(-J) = (-1)^{\dim V} \det(J)$, which implies $(-1)^{\dim V} = 1$ for non-zero determinant.}.
    \item \textbf{Euler Characteristic $\chi(M)$}: Defined as the alternating sum of Betti numbers: $\chi(M) = \sum_{k=0}^n (-1)^k b_k$.
\end{itemize}

\textbf{Geometric Intuition:}
The Euler characteristic is a topological invariant that counts the "alternating number of holes" in a space. For a closed manifold, Poincaré Duality forces a symmetry in this count. In our 6D case, $b_0=b_6$, $b_1=b_5$, and $b_2=b_4$. The "odd one out" is the middle Betti number $b_3$. 
The core of the problem lies in the nature of the "intersections" of 3-dimensional cycles. In a 6D manifold, the intersection of two 3-cycles is a 0-dimensional object (points). Because the dimension is odd ($3$), the order in which we intersect them matters: $A \cap B = -B \cap A$. This anti-symmetry is characteristic of a symplectic structure. Just as a symplectic vector space must be even-dimensional, the space of 3-cycles (measured by $b_3$) must have an even rank. This parity constraint on $b_3$ is what ultimately determines the possible values of $\chi(M)$.

\textbf{Formal Proof:}

\textbf{1. Determining Betti Numbers via Duality and Constraints:}
Since $M$ is connected and closed, the top and bottom Betti numbers are fixed:
\sidenote{$b_0 = 1$ and $b_6 = 1$.}
We are given that $M$ is simply connected, which implies $H_1(M; \mathbb{Z}) = 0$. Thus, $b_1 = 0$. By Poincaré Duality:
\sidenote{$b_5 = b_1 = 0$.}
Next, we are given $H_2(M; \mathbb{Z}) = \mathbb{Z}_2$. The Betti number $b_2$ is the rank of the free part of $H_2$. Since $\mathbb{Z}_2$ is purely torsion:
\sidenote{$b_2 = \text{rank}(\mathbb{Z}_2) = 0$.}
Again, by Poincaré Duality, the Betti number of the dual dimension must match:
\sidenote{$b_4 = b_2 = 0$.}

\textbf{2. Evaluating the Euler Characteristic:}
The formula for the Euler characteristic of a 6-manifold is:
\[
\chi(M) = b_0 - b_1 + b_2 - b_3 + b_4 - b_5 + b_6.
\]
Substituting the known values:
\sidenote{$\chi(M) = 1 - 0 + 0 - b_3 + 0 - 0 + 1 = 2 - b_3$.}

\textbf{3. The Symplectic Nature of the Middle Dimension:}
Consider the cohomology group with real coefficients $H^3(M; \mathbb{R})$. The cup product defines a bilinear form:
\[
Q: H^3(M; \mathbb{R}) \times H^3(M; \mathbb{R}) \to \mathbb{R}, \quad Q(\alpha, \beta) = \langle \alpha \cup \beta, [M] \rangle.
\]
\sidenote{Construction of the Symplectic Form:
For $\alpha, \beta \in H^3(M; \mathbb{R})$, the form $Q(\alpha, \beta)$ is defined by evaluating the cup product $\alpha \cup \beta \in H^6(M; \mathbb{R})$ on the fundamental class $[M] \in H_6(M; \mathbb{R})$. 
In the de Rham framework, this corresponds to $Q([\omega], [\eta]) = \int_M \omega \wedge \eta$.
Graded commutativity implies $\alpha \cup \beta = (-1)^{3 \cdot 3} \beta \cup \alpha = -\beta \cup \alpha$, establishing skew-symmetry. 
Non-degeneracy follows from Poincaré Duality, which asserts that the map $H^3(M; \mathbb{R}) \to \text{Hom}(H^3(M; \mathbb{R}), \mathbb{R})$ given by $\alpha \mapsto Q(\alpha, \cdot)$ is an isomorphism.}
Thus, $Q$ is a \textbf{skew-symmetric} bilinear form. Poincaré Duality further guarantees that $Q$ is \textbf{non-degenerate}.
\footnote{A non-degenerate skew-symmetric form is called a symplectic form. A standard result in linear algebra states that any real vector space admitting a non-degenerate skew-symmetric bilinear form must have an even dimension. This is because such a space can be decomposed into pairs of hyperbolic planes.}
This implies that the dimension of $H^3(M; \mathbb{R})$, which is exactly the Betti number $b_3$, must be even.
\sidenote{$b_3 = 2k$ for some integer $k \ge 0$.}

\textbf{4. Conclusion:}
Plugging $b_3 = 2k$ back into our expression for the Euler characteristic:
\sidenote{$\chi(M) = 2 - 2k = 2(1-k)$.}
This shows that $\chi(M)$ is an \textbf{even integer}.
Since $-1$ is an odd integer, it is impossible for $\chi(M)$ to equal $-1$.
\end{solution}

\begin{remark}[Step-by-Step Cohomology of Standard Manifolds]
To build intuition for Betti numbers and Poincaré Duality, let's explicitly calculate the groups and Euler characteristics for standard spaces:
\begin{enumerate}
    \item \textbf{Spheres $S^n$}:
    The cohomology groups are $H^0 \cong \mathbb{Z}$ and $H^n \cong \mathbb{Z}$.
    \begin{itemize}
        \item For $S^6$: Betti numbers are $b_0=1, b_6=1$, and $b_k=0$ for $0 < k < 6$.
        \item Euler characteristic: $\chi(S^6) = 1 - 0 + 0 - 0 + 0 - 0 + 1 = 2$.
        \item General rule: $\chi(S^n) = 1 + (-1)^n$, which is 2 for $n$ even and 0 for $n$ odd.
    \end{itemize}
    
    \item \textbf{Complex Projective Space $\mathbb{CP}^n$}:
    The cohomology exists only in even dimensions: $H^{2k} \cong \mathbb{Z}$ for $k=0, \dots, n$.
    \begin{itemize}
        \item For $\mathbb{CP}^3$ (dim 6): Betti numbers are $b_0=1, b_2=1, b_4=1, b_6=1$, and $b_k=0$ otherwise.
        \item Euler characteristic: $\chi(\mathbb{CP}^3) = 1 - 0 + 1 - 0 + 1 - 0 + 1 = 4$.
        \item General rule: $\chi(\mathbb{CP}^n) = n+1$.
    \end{itemize}
    
    \item \textbf{Torus $T^n = (S^1)^n$}:
    The Betti numbers are given by the binomial coefficients $\binom{n}{k}$.
    \begin{itemize}
        \item For $T^3$ (dim 3): $b_0 = \binom{3}{0}=1, b_1 = \binom{3}{1}=3, b_2 = \binom{3}{2}=3, b_3 = \binom{3}{3}=1$.
        \item Euler characteristic: $\chi(T^3) = 1 - 3 + 3 - 1 = 0$.
        \item For $T^6$: The middle Betti number is $b_3 = \binom{6}{3} = 20$.
        \item General rule: $\chi(T^n) = 0$ for all $n \ge 1$.
    \end{itemize}
    
    \item \textbf{Product of Spheres $S^3 \times S^3$}:
    We use the Künneth formula: $b_k(X \times Y) = \sum_{i+j=k} b_i(X)b_j(Y)$.
    \begin{itemize}
        \item For $S^3 \times S^3$ (dim 6):
        \begin{itemize}
            \item $b_0 = b_0(S^3)b_0(S^3) = 1 \cdot 1 = 1$.
            \item $b_3 = b_0(S^3)b_3(S^3) + b_3(S^3)b_0(S^3) = 1 \cdot 1 + 1 \cdot 1 = 2$.
            \item $b_6 = b_3(S^3)b_3(S^3) = 1 \cdot 1 = 1$.
        \end{itemize}
        \item Euler characteristic: $\chi(S^3 \times S^3) = 1 - 0 + 0 - 2 + 0 - 0 + 1 = 0$.
    \end{itemize}
\end{enumerate}
\end{remark}

\begin{exercise}
\label{geometry:2024:4}
Let $(M, g)$ be a closed oriented $n$-dimensional Riemannian manifold. Let $p \in M$ and $\mathrm{Ric}_p$ be the Ricci curvature tensor at $p$, $S_p$ be the scalar curvature at $p$ which is defined to be $S_p := \frac{1}{n} \mathrm{Tr}_g(\mathrm{Ric}_p)$. Prove that the scalar curvature $S_p$ at $p \in M$ is given by
\[
S_p = \frac{1}{\omega_{n-1}} \int_{S^{n-1}} \mathrm{Ric}_p(V, V) d S^{n-1},
\]
where $\omega_{n-1}$ is the area of the unit sphere $S^{n-1}$ in $T_p M$, $V \in S^{n-1}$ are unit vector fields, and $d S^{n-1}$ is the area element on $S^{n-1}$.
\end{exercise}

\begin{solution}
\textbf{Prerequisites \& Definitions:}
\begin{itemize}
    \item \textbf{Ricci Curvature $\mathrm{Ric}_p$}: This is a symmetric $(0,2)$-tensor obtained by taking a trace of the Riemann curvature tensor $R$. Formally, for tangent vectors $X, Y$, the Ricci tensor is defined as:
    \[ \mathrm{Ric}(X, Y) = \text{Tr}(Z \mapsto R(Z, X)Y) \]
    In local coordinates, this corresponds to the contraction:
    \[ R_{ij} = R^k_{ikj}, \quad \text{where } R^k_{lij} \text{ are components of the Riemann tensor.} \]
    \footnote{Geometrically, $\mathrm{Ric}(V, V)$ is the sum of sectional curvatures of $(n-1)$ independent planes containing $V$. It measures the rate at which the volume of a narrow cone of geodesics starting at $p$ in direction $V$ deviates from the Euclidean case.}
    \item \textbf{Scalar Curvature $S_p$}: This is the simplest curvature invariant, obtained by taking the trace of the Ricci tensor. It measures how the volume of a small ball in the manifold differs from a ball in Euclidean space\footnote{Specifically, $\frac{\text{Vol}(B_{\epsilon}(p))}{\text{Vol}_{\text{Eucl}}(B_{\epsilon})} = 1 - \frac{S}{6(n+2)}\epsilon^2 + O(\epsilon^4)$.}.
    \item \textbf{Symmetry of Sphere Integrals}: When integrating components of a vector $V = (v_1, \dots, v_n)$ over the unit sphere, the integral $\int_{S^{n-1}} v_i v_j dS$ is zero if $i \neq j$ because the sphere is symmetric under reflections $v_i \to -v_i$.
\end{itemize}

\textbf{Geometric Intuition:}
The Ricci curvature $\mathrm{Ric}_p(V, V)$ tells you how the manifold is curving in the "direction" $V$. The scalar curvature is the total curvature at point $p$, which we can think of as the "average" of all these directional curvatures. By integrating the Ricci curvature over all possible unit directions (the sphere $S^{n-1}$) and dividing by the total area $\omega_{n-1}$, we are literally computing the average value of the Ricci curvature.

\textbf{Formal Proof:}

Let $p \in M$. In the tangent space $T_p M$, we can choose an orthonormal basis $\{e_1, e_2, \dots, e_n\}$ that diagonalizes the Ricci tensor:
\[ \mathrm{Ric}_p(e_i, e_j) = R_{ii} \delta_{ij} \]
Any unit vector $V \in S^{n-1} \subset T_p M$ can be written in terms of its components in this basis:
\[ V = \sum_{i=1}^n v_i e_i, \quad \text{with} \quad \sum_{i=1}^n v_i^2 = 1 \]
Now, let's expand the expression for the Ricci curvature in the direction $V$:
\[ \mathrm{Ric}_p(V, V) = \sum_{i,j} \mathrm{Ric}_p(v_i e_i, v_j e_j) = \sum_{i,j} R_{ij} v_i v_j \]
Since the basis is orthonormal and diagonalizes the Ricci tensor, $R_{ij} = 0$ for $i \neq j$:
\[ \mathrm{Ric}_p(V, V) = \sum_{i=1}^n R_{ii} v_i^2 \]
We wish to compute the integral of this expression over the unit sphere $S^{n-1}$:
\[ \int_{S^{n-1}} \mathrm{Ric}_p(V, V) dS = \sum_{i=1}^n R_{ii} \int_{S^{n-1}} v_i^2 dS \]
By the symmetry of the sphere, the integral of $v_i^2$ is the same for all $i$:
\[ \sum_{k=1}^n \int_{S^{n-1}} v_k^2 dS = \int_{S^{n-1}} \left(\sum v_k^2\right) dS = \int_{S^{n-1}} 1 dS = \omega_{n-1} \]
Since there are $n$ such identical integrals, we have:
\[ \int_{S^{n-1}} v_i^2 dS = \frac{\omega_{n-1}}{n} \]
Substituting this back into the integral of the Ricci curvature:
\[ \int_{S^{n-1}} \mathrm{Ric}_p(V, V) dS = \sum_{i=1}^n R_{ii} \left(\frac{\omega_{n-1}}{n}\right) = \frac{\omega_{n-1}}{n} \sum_{i=1}^n R_{ii} \]
The sum $\sum R_{ii}$ is the trace of the Ricci tensor:
\[ \sum_{i=1}^n R_{ii} = \mathrm{Tr}_g(\mathrm{Ric}_p) \]
By the definition given in the problem, $S_p = \frac{1}{n} \mathrm{Tr}_g(\mathrm{Ric}_p)$:
\[ \int_{S^{n-1}} \mathrm{Ric}_p(V, V) dS = \omega_{n-1} \cdot S_p \]
Rearranging the terms, we obtain the final formula:
\[ S_p = \frac{1}{\omega_{n-1}} \int_{S^{n-1}} \mathrm{Ric}_p(V, V) dS^{n-1} \]
\end{solution}

\begin{exercise}
\label{geometry:2024:5}
Let $S^n$ be the $n$-dimensional sphere with $n \ge 2$, and let $G$ be a finite group that acts freely on $S^n$. Suppose $G$ is non-trivial. Then,
(i) Compute the homotopy groups of the quotient space $\pi_i(S^n/G)$ for $0 \le i \le n$.
(ii) Suppose $n$ is even. Prove that $G$ is isomorphic to $\mathbb{Z}_2$.
(iii) Suppose $n$ is odd. Show that $G$ cannot be isomorphic to $\mathbb{Z}_p \times \mathbb{Z}_p$ for $p$ a prime number.
\end{exercise}

\begin{solution}
\textbf{Prerequisites \& Definitions:}
\begin{itemize}
    \item \textbf{Free Group Action}: A group $G$ acts freely on a space $X$ if the only element that has a fixed point is the identity element $e$\footnote{Formally, for all $g \in G \setminus \{e\}$ and all $x \in X$, $g \cdot x \neq x$.}.
    \item \textbf{Quotient Space $X/G$}: This is the set of orbits of the action.
    \sidenote{Definition of the Quotient:
    The space $X/G$ is defined by the equivalence relation $x \sim y$ if there exists $g \in G$ such that $y = g \cdot x$. The set of all such points $\{g \cdot x \mid g \in G\}$ is called the \textbf{orbit} of $x$. 
    The quotient space is equipped with the \textbf{quotient topology}, where a set $U \subset X/G$ is open if and only if its preimage $p^{-1}(U) \subset X$ is open.
    Because $G$ is finite and acts freely and continuously, the projection $p: X \to X/G$ is an open map and a local homeomorphism, ensuring $X/G$ inherits the manifold structure and Hausdorff property of $X$.}
    \item \textbf{Covering Space}: If a finite group $G$ acts freely on a manifold $X$, the projection $p: X \to X/G$ is a regular covering map\footnote{The quotient space $X/G$ inherits the local manifold structure of $X$, and the fiber over any point in the quotient consists of $|G|$ points in $X$.}.
    \item \textbf{Homotopy Groups $\pi_i$}: For a covering $p: \tilde{X} \to X$ with $\tilde{X}$ simply connected, $\pi_1(X) \cong G$ (the deck transformation group) and $\pi_i(X) \cong \pi_i(\tilde{X})$ for $i \ge 2$\footnote{This follows from the long exact sequence of homotopy groups for a covering (which is a fibration with discrete fiber).}.
\end{itemize}

\textbf{Geometric Intuition:}
Imagine the sphere $S^n$. A free action means $G$ moves points around like a "shuffling" operation. For $n$ even, the Hairy Ball Theorem tells us we can't have a non-vanishing vector field, which is a hint that it's hard to "rotate" the sphere without fixing some points (like the poles). The only way to move every point without fixing any is a "flip" like the antipodal map $x \to -x$. Since doing this flip twice gets you back to where you started, the group can only have two elements: the identity and the flip. For $n$ odd, the constraints are looser (you can rotate an odd sphere), but the group still can't be too "complex" (like a grid $\mathbb{Z}_p^2$) because it has to "fit" into the homology of a sphere-like space.

\textbf{Formal Proof:}

\textbf{(i) Step-by-Step Derivation of the Homotopy Groups of $S^n/G$:}

\textbf{1. The Covering Space Structure:}
Since $G$ is a finite group acting freely on $S^n$, the projection map $p: S^n \to S^n/G$ is a \textbf{regular covering map}. 
For $n \ge 2$, the sphere $S^n$ is simply connected ($\pi_1(S^n) = 0$). 
\sidenote{A covering space with a simply connected total space is called the \textbf{universal cover}. This is the most "unrolled" version of the space.}
Thus, $p: S^n \to S^n/G$ is the universal covering of the quotient space $X = S^n/G$.

\textbf{2. Using the Long Exact Sequence of a Fibration:}
Every covering map is a \textbf{fibration} where the fiber $F$ is a discrete set of points (the points in an orbit). For a regular covering, the fiber $F$ is in one-to-one correspondence with the group $G$.
The homotopy groups of any fibration $F \to E \to B$ are linked by a long exact sequence:
\[ \dots \to \pi_i(F) \to \pi_i(E) \to \pi_i(B) \xrightarrow{\partial} \pi_{i-1}(F) \to \dots \]
In our case, $F = G$, $E = S^n$, and $B = S^n/G$.

\textbf{3. Higher Homotopy Groups ($i \ge 2$):}
Because $G$ is a discrete group, its points are isolated. A "continuous map from a sphere" into a discrete set must be constant.
\sidenote{Discrete Fiber Homotopy: For any discrete space $F$, $\pi_i(F) = 0$ for all $i \ge 1$.}
Substituting this into the sequence for $i \ge 2$:
\[ \dots \to 0 \to \pi_i(S^n) \to \pi_i(S^n/G) \to 0 \to \dots \]
By the property of \textbf{exactness} (the image of one map is the kernel of the next), a map between two groups surrounded by zeros must be an \textbf{isomorphism}.
\sidenote{Result for $i \ge 2$: $\pi_i(S^n/G) \cong \pi_i(S^n)$.}
Using the known homotopy groups of the sphere:
\begin{itemize}
    \item For $2 \le i < n$, $\pi_i(S^n/G) \cong \pi_i(S^n) = 0$.
    \item For $i = n$, $\pi_n(S^n/G) \cong \pi_n(S^n) \cong \mathbb{Z}$.
\end{itemize}

\textbf{4. The Fundamental Group ($i = 1$):}
For the first dimension, the end of the long exact sequence is:
\[ \pi_1(S^n) \to \pi_1(S^n/G) \xrightarrow{\partial} \pi_0(G) \to \pi_0(S^n) \]
We substitute the values for the sphere:
\begin{itemize}
    \item $\pi_1(S^n) = 0$ (simply connected).
    \item $\pi_0(S^n) = 0$ (path-connected, so there is only one "component").
    \item $\pi_0(G) = G$ (each element of the group is a separate component).
\end{itemize}
The sequence becomes:
\[ 0 \to \pi_1(S^n/G) \xrightarrow{\cong} G \to 0 \]
\sidenote{Conclusion for $i=1$: $\pi_1(S^n/G) \cong G$.}
Intuitively, since the top space has no "holes" ($\pi_1=0$), all the "loops" in the quotient space must come from moving from one point in an orbit to another, which is exactly what the group $G$ does.

\textbf{5. The 0-th Homotopy Group ($i = 0$):}
Since $S^n$ is path-connected and the map $p$ is surjective, the quotient $S^n/G$ is also path-connected.
\sidenote{Result for $i = 0$: $\pi_0(S^n/G) = 0$.}

\textbf{(ii) Case $n$ is even:}
Let $g \in G$ be any non-identity element. Since the action is free, $g$ has no fixed points.
According to the \textbf{Lefschetz Fixed Point Theorem}:
\footnote{The Lefschetz number $L(g)$ of a map $g$ is defined as $\sum_{i=0}^n (-1)^i \text{Tr}(g_* | H_i(X; \mathbb{Q}))$. If $g$ has no fixed points, then $L(g) = 0$.}
For the sphere $S^n$, the only non-zero homology groups are $H_0 \cong \mathbb{Q}$ and $H_n \cong \mathbb{Q}$.
\sidenote{$L(g) = \text{Tr}(g_* | H_0) + (-1)^n \text{Tr}(g_* | H_n) = 1 + (-1)^n \deg(g)$.}
For $n$ even, we have $L(g) = 1 + \deg(g)$.
Since $g$ is a homeomorphism, $\deg(g) = \pm 1$.
If $g$ has no fixed points, $L(g)$ must be $0$, which forces $\deg(g) = -1$.
Now, consider any $g, h \in G \setminus \{e\}$. Then $\deg(g) = -1$ and $\deg(h) = -1$.
The degree of the composition is:
\sidenote{$\deg(gh) = \deg(g)\deg(h) = (-1)(-1) = 1$.}
If $gh \neq e$, then $gh$ would have a non-zero Lefschetz number $L(gh) = 1+1 = 2$, meaning it must have a fixed point.
This contradicts the free action. Thus, we must have $gh = e$ for all non-identity $g, h$.
This implies every non-identity element is its own inverse, and there is only one such element.
Thus, $G \cong \mathbb{Z}_2$.

\textbf{(iii) Case $n$ is odd:}
Suppose $G \cong \mathbb{Z}_p \times \mathbb{Z}_p$ acts freely on $S^n$.
A finite group $G$ acting freely on a sphere $S^n$ must satisfy the \textbf{p-periodicity} condition:
\footnote{A group $G$ has $p$-periodic cohomology if its cohomology with coefficients in $\mathbb{Z}_p$ is periodic. Smith's theorem states that a finite group acting freely on a sphere $S^n$ cannot have any subgroup isomorphic to $\mathbb{Z}_p \times \mathbb{Z}_p$.}
The cohomology ring $H^*(\mathbb{Z}_p \times \mathbb{Z}_p; \mathbb{Z}_p)$ is quite large; its dimensions grow linearly with the degree.
In contrast, if $G$ acts freely on $S^n$, the quotient $S^n/G$ is a manifold whose cohomology must vanish above dimension $n$.
Using the spectral sequence for the covering, we find that the cohomology of $G$ must be periodic with period $n+1$\footnote{Specifically, if $X$ is a homotopy sphere, $H^k(G; \mathbb{Z}) \cong H^k(X/G; \mathbb{Z})$ for $0 < k < n$. For $k > n$, periodicity kicks in. $\mathbb{Z}_p^2$ does not have this property.}.
This periodicity is impossible for $\mathbb{Z}_p \times \mathbb{Z}_p$.
Thus, $G$ cannot be isomorphic to $\mathbb{Z}_p \times \mathbb{Z}_p$.
\end{solution}

\begin{exercise}
\label{geometry:2024:6}
Let $M$ be a closed oriented Riemannian manifold, where $g_t$ is a family of smooth Riemannian metrics smoothly depending on $t \in (-\epsilon, \epsilon)$. Suppose there exists a family of eigenfunctions $f_t$ and eigenvalues $\lambda_t$ smoothly depending on $t$ such that
\[
\Delta_{g_t} f_t = \lambda_t f_t,
\]
where $\Delta_{g_t}$ is the Laplace-Beltrami operator defined using the Riemannian metric $g_t$. Additionally, assume that $f_0$ is not a constant function. We define $\dot{\lambda} := \left.\frac{d}{dt}\right|_{t=0} \lambda_t$ and $\dot{\Delta} := \left.\frac{d}{dt}\right|_{t=0} \Delta_{g_t}$. Prove the following:
(i) As $\lambda_0$ is an eigenvalue of $\Delta_{g_0}$, let $V_{\lambda_0} := \mathrm{Ker}(\Delta_{g_0} - \lambda_0)$ be the eigenspace of $\lambda_0$, and let $\Pi : L^2(M, g_0) \to V_{\lambda_0}$ be the orthogonal projection onto the eigenspace. Prove that $\dot{\lambda}$ is an eigenvalue of the operator $\Pi \circ \dot{\Delta} : V_{\lambda_0} \to V_{\lambda_0}$.
(ii) Let $\varphi_t : M \to M$ be a 1-parameter family of diffeomorphisms of $M$ and assume $g_t = \varphi_t^* g_0$. Prove that $\dot{\lambda} = 0$.
\end{exercise}

\begin{solution}
\textbf{Prerequisites \& Definitions:}
\begin{itemize}
    \item \textbf{Laplace-Beltrami Operator $\Delta_g$}: This is the generalization of the Laplacian to a Riemannian manifold. Its eigenvalues $\lambda$ are always non-negative, and its eigenfunctions $f$ represent the "modes of vibration" of the manifold\footnote{For a function $u$, $\Delta u = -\text{div}(\text{grad } u)$. On a closed manifold, the spectrum is discrete: $0 = \lambda_0 < \lambda_1 \le \lambda_2 \dots \to \infty$.}.
    \item \textbf{Eigenvalue Variation}: When the metric $g_t$ changes, the operator $\Delta_{g_t}$ changes, and consequently its eigenvalues $\lambda_t$ and eigenfunctions $f_t$ also vary. The study of this variation is central to spectral geometry.
    \item \textbf{Orthogonal Projection $\Pi$}: In the Hilbert space $L^2(M, g_0)$, $\Pi$ maps any function to its "shadow" on the subspace $V_{\lambda_0}$\footnote{If $\{u_1, \dots, u_k\}$ is an orthonormal basis for $V_{\lambda_0}$, then $\Pi(h) = \sum_{j=1}^k \langle h, u_j \rangle u_j$.}.
\end{itemize}

\textbf{Geometric Intuition:}
Imagine the manifold is a bell. The eigenvalues are the frequencies you hear when you strike it. If you slightly change the shape of the bell (change the metric $g_t$), the frequencies will change. Part (i) tells us exactly how much the frequency shifts based on the "speed" at which the shape is changing ($\dot{\Delta}$). Part (ii) says that if the "change" is just a relabeling of the points on the bell (a diffeomorphism $\varphi_t$) without actually changing its physical geometry, the frequencies won't change at all—a rotated bell sounds exactly the same as the original.

\textbf{Formal Proof:}

\textbf{(i) The First-Order Variation Formula:}
We start with the eigenvalue equation at each time $t$:
\sidenote{$\Delta_{g_t} f_t = \lambda_t f_t$.}
Differentiating this equation with respect to $t$ at $t=0$ using the product rule:
\sidenote{$\dot{\Delta} f_0 + \Delta_{g_0} \dot{f} = \dot{\lambda} f_0 + \lambda_0 \dot{f}$.}
Rearranging the terms to group the variation of the operator and the variation of the eigenvalue:
\sidenote{$(\dot{\Delta} - \dot{\lambda}) f_0 = (\lambda_0 - \Delta_{g_0}) \dot{f}$.}
Now, we apply the orthogonal projection $\Pi$ onto the eigenspace $V_{\lambda_0}$ to both sides of the equation.
\sidenote{$\Pi(\dot{\Delta} f_0) - \Pi(\dot{\lambda} f_0) = \Pi((\lambda_0 - \Delta_{g_0}) \dot{f})$.}
Since $f_0 \in V_{\lambda_0}$, we have $\Pi(f_0) = f_0$ and $\Pi(\dot{\lambda} f_0) = \dot{\lambda} f_0$.
For the right-hand side, consider the fact that $\Delta_{g_0}$ is self-adjoint on $L^2(M, g_0)$.
\footnote{For any $u \in V_{\lambda_0}$ and any $v \in L^2$, we have $\langle (\lambda_0 - \Delta_{g_0}) v, u \rangle = \langle v, (\lambda_0 - \Delta_{g_0}) u \rangle = \langle v, 0 \rangle = 0$. This means the range of $(\lambda_0 - \Delta_{g_0})$ is orthogonal to $V_{\lambda_0}$.}
Therefore, the projection of any element in the image of $(\lambda_0 - \Delta_{g_0})$ onto $V_{\lambda_0}$ is zero.
\sidenote{$\Pi((\lambda_0 - \Delta_{g_0}) \dot{f}) = 0$.}
The equation simplifies to:
\sidenote{$\Pi(\dot{\Delta} f_0) = \dot{\lambda} f_0$.}
This demonstrates that $f_0$ is an eigenfunction of the operator $\Pi \circ \dot{\Delta}$ with eigenvalue $\dot{\lambda}$.

\textbf{(ii) Invariance Under Diffeomorphisms:}
If $g_t = \varphi_t^* g_0$ for some diffeomorphisms $\varphi_t$, then by the naturality of the Laplace-Beltrami operator, we have:
\sidenote{$\Delta_{g_t} = \varphi_t^* \circ \Delta_{g_0} \circ (\varphi_t^*)^{-1}$.}
This means that for every $t$, the operator $\Delta_{g_t}$ is \textbf{isometrically conjugate} to $\Delta_{g_0}$.
\footnote{Two conjugate operators have the same spectrum. In this case, $\Delta_{g_t} (\varphi_t^* f) = \lambda (\varphi_t^* f)$ if and only if $\Delta_{g_0} f = \lambda f$.}
Therefore, the eigenvalues of $\Delta_{g_t}$ are identical to the eigenvalues of $\Delta_{g_0}$ for all $t \in (-\epsilon, \epsilon)$.
\sidenote{$\lambda_t = \lambda_0$ for all $t$.}
It follows immediately that the derivative of the eigenvalue with respect to time must be zero:
\sidenote{$\dot{\lambda} = \frac{d}{dt} \lambda_t = 0$.}
\end{solution}
